\cotefile{156-3}
\scanpage{10}\mspage{9}
% ink top, lines 1–3
un \uline{demi}-\uline{intervalle} s'il admet un plus petit \unsure{élément} mais pas un plus grand, ou bien l'inverse.

% ink upper middle, NB, lines 1–4
\deleted{\unsure{L'opposé|L'oppo}}\inserted{\uline{NB}} L'opposé d'un tronçon, segment, intervalle, demi-intervalle est itou. L'opposé d'un $1/2$ intervalle fermé à g. est fermé à droite, et \unsure{inversement}.
% ink left, slanting upwards beside NB
% slanted
\margin{Dans \unsure{ma|sa} \unsure{actuelle|initiale} \unsure{terminologie}, tronçon \unsure{est|être|était|etre} \unsure{``ens|ens|un|équiv} \uline{\unsure{biordre|biordonné}}}

% ink middle, Tout, lines 1–3
\deleted{\unsure{Lorsque|Lorsqu.}}\inserted{Tout} tronçon est l'un des trois (segment, $1/2$ intervalle, intervalle) à l'exclusion \unsure{des} deux autres.

% ink lower middle, Si L
Si $\mathcal L$ est un tronçon, on pose $\deleted{\unsure{\partial L}=}$
\[
\begin{aligned}
\unsure{\partial L|\partial}&=\bigl\{\,x\in L\mid
\begin{aligned}&x\text{ est plus petit élément}\\
&\text{ou plus grand \unsure{élément|élt}}\end{aligned}\,\bigr\}
\end{aligned}
\]
\deleted{\illegible{1}}

% ink bottom, Soit and first equality
Soit $\mathcal L$ un tronçon, et
\[
\begin{aligned}
\unsure{\mathscr S|S}_0(L)=\mathrm{Drap}_2(L)&=\{\,\unsure{E|\mathcal E}\in\unsure{\mathrm P|\mathscr P}_2(L)\mid E\text{ tot ordonné}\,\}\\
&=\{\,\{x,y\}\mid x<y\,\}
\end{aligned}
\]
% ink bottom, final equality
\[
\text{D'où}\qquad \unsure{\mathrm{Modf}|\mathrm{Mod}_{f}}(L)=\mathrm{Drap}(L)
=\{\,S\subset\mathrm P_f(L)\mid S\,\substack{\text{tot}\\\text{ord}}\,\}
\]
